\documentclass[10pt,a4paper]{article}
\usepackage{fontspec,amsmath,amssymb,graphicx}
\usepackage[margin=18mm]{geometry}
\setmainfont{Times New Roman}
\setlength{\parindent}{0pt}
\setlength{\parskip}{5pt}
\allowdisplaybreaks
\sloppy
\emergencystretch=2em
\begin{document}
\section*{BT\_1\_2\_3\_05 — Lời giải bổ sung}
\subsection*{Kiến thức và trình tự giải}
\textbf{Mục tiêu.} Đưa phương trình về ODE bằng cách giữ cố định các biến không được lấy đạo hàm, hoặc đặt một đạo hàm/hợp thức trung gian.

\textbf{Bước 1 — Chọn biến và tham số.} Với $u_x=f(x,y)$, cố định $y$ và tích phân theo $x$. Khi đó “hằng số” tích phân có thể khác nhau ở mỗi $y$, nên phải viết $C(y)$. Với ba biến, tích phân theo $x$ cho một hàm $C(y,z)$.

\textbf{Bước 2 — Tích phân đúng số lần.} Chẳng hạn $u_{xx}=f$ cho $u_x=\int f\,dx+A(y)$, rồi $u=\int(\int f\,dx)dx+xA(y)+B(y)$. Với đạo hàm hỗn hợp, có thể đặt $v=u_x$ hoặc $v=yu_y+2u$ nếu PDE chứa đúng $v_x$. Mỗi lần tích phân phải giữ hàm tùy ý theo các biến còn lại.

\textbf{Bước 3 — Giải ODE tuyến tính nếu còn hạng chứa ẩn.} Với $u_x+P(x)u=Q(x,y)$, đặt $M(x)=e^{\int P(x)dx}$. Do $M'=PM$, quy tắc đạo hàm tích cho

\[
(Mu)_x=M(u_x+Pu)=MQ,\qquad
u=M^{-1}\left[\int MQ\,dx+C(y)\right].
\]

Ghi rõ $P$, $M$ và cách tính tích phân trong từng bài. Với ODE cấp 2 có hệ số hằng theo biến đang tích phân, thử $e^{\lambda x}$ và giải phương trình đặc trưng; nghiệm kép cần thêm $xe^{\lambda x}$.

\textbf{Bước 4 — Xét trường hợp suy biến trước khi chia.} Chia cho $x$, $y$ hoặc $u$ chỉ hợp lệ ở nơi chúng khác 0. Trên đường bị loại, dùng phương trình gốc để tìm ràng buộc. Nếu cần nghiệm trơn qua đường đó, xét giới hạn và đạo hàm của công thức vừa tìm.

\textbf{Bước 5 — Dùng dữ liệu và kiểm tra.} Thế điều kiện đầu/biên để xác định các hàm tùy ý, rồi lấy đạo hàm của nghiệm và thay lại phương trình. Nêu rõ miền và độ trơn cần thiết. Nếu chỉ tìm nghiệm dạng mũ, đó là một họ nghiệm thử; không tự gọi là toàn bộ nghiệm tổng quát.

\textbf{Bổ sung cho ý tổng bình phương đạo hàm.} Với nghiệm thực, $u_x^2+u_y^2=0$ buộc từng đạo hàm bằng 0, nên $u$ hằng trên mỗi thành phần liên thông. Nếu vế phải là $k<0$, vô nghiệm thực; nếu $k>0$, đây là eikonal $|\nabla u|^2=k$. Nghiệm affine $u=ax+by+c$ cần $a^2+b^2=k$, nhưng chỉ là một họ. Muốn dựng nghiệm cục bộ từ dữ liệu trên đường $\gamma(s)$, chọn gradient ban đầu $P(s)$ thỏa $|P|^2=k$, $P\cdot\gamma'=h'$. Đặc trưng cho $X=\gamma(s)+2\tau P(s)$, $U=h(s)+2k\tau$; chỉ khử được tham số nơi Jacobian khác 0. Cách này được áp dụng riêng ở ý (d).
\subsection*{Lời giải áp dụng cho bài}
(a) Xem $y$ là tham số: $u_x=-u$. Nhân $e^x$ cho $(e^xu)_x=0$, nên $u=e^{-x}C(y)$. Kiểm tra $u_x+u=0$.

(b) Xem $x$ là tham số. ODE theo $y$ là $u_{yy}-x^2u=0$, có phương trình đặc trưng $\lambda^2-x^2=0$. Với $x\ne0$, $u=A(x)e^{xy}+B(x)e^{-xy}$. Với $x=0$, ODE thành $u_{yy}=0$, nên vết là affine theo $y$. Một dạng thống nhất, không mất nghiệm ở $x=0$, là
\[u=A(x)\cosh(xy)+B(x)\frac{\sinh(xy)}x,\]
trong đó $\sinh(xy)/x=y$ tại $x=0$ theo giới hạn. Hệ số $A,B$ đủ trơn cho lớp nghiệm mong muốn.

(c) Đặt $v=u_x$, được $v_x-v=0$, nên $v=B(y)e^x$. Tích phân theo $x$ cho $u=A(y)+B(y)e^x$. Đạo hàm lại cho $u_{xx}=u_x$.

(d) Với nghiệm thực, $u_x^2+u_y^2=0$ là tổng hai số không âm; buộc $u_x=u_y=0$. Do đó $u$ hằng trên từng thành phần liên thông. Nếu vế phải là $k\ne0$: $k<0$ thì vô nghiệm thực; $k>0$ là eikonal $|\nabla u|=\sqrt k$. Một họ nghiệm đơn giản $u=ax+by+c$, $a^2+b^2=k$, nhưng không phải toàn bộ nghiệm.

Để dựng nghiệm cục bộ tổng quát từ dữ liệu $u(\gamma(s))=h(s)$, chọn $P(s)$ thỏa $|P(s)|^2=k$ và $P(s)\cdot\gamma'(s)=h'(s)$. Hệ đặc trưng của $F=p^2+q^2-k$ có $P'=0$, $X'=2P$, $U'=2k$, nên
\[X=\gamma(s)+2tP(s),\qquad U=h(s)+2kt.\]
Khử $(s,t)$ ở nơi Jacobian của $X$ khác 0, được nghiệm $u(X)=U$. Đây là dạng áp dụng phương pháp đặc trưng, phân biệt với chỉ đưa ra ví dụ affine.

\textbf{Chi tiết xử lý điểm $x=0$ ở (b).} Chọn hai hệ số từ các vết tại $y=0$: $A(x)=u(x,0)$, $B(x)=u_y(x,0)$. Khi $x\ne0$, giải hệ $C+D=A$, $x(C-D)=B$ cho $C=(A+B/x)/2,D=(A-B/x)/2$, rồi gom hai hàm mũ thành $A\cosh(xy)+B\sinh(xy)/x$. Khai triển $\sinh(xy)/x=y+x^2y^3/3!+\cdots$ cho giới hạn $y$ và vết $u(0,y)=A(0)+B(0)y$; không loại nghiệm tuyến tính này khi hai nghiệm mũ trùng nhau.\end{document}